% !TEX spellcheck = en_US
% !TEX spellcheck = LaTeX
\documentclass[letterpaper,10pt,english]{article}
%\usepackage{%
amsfonts,%
amsmath,%
amssymb,%
amsthm,%
babel,%
bbm,%
%biblatex,%
caption,%
centernot,%
color,%
enumerate,%
%enumitem,%
epsfig,%
epstopdf,%
etex,%
fancybox,%
framed,%
fullpage,%
%geometry,%
graphicx,%
hyperref,%
latexsym,%
mathptmx,%
mathtools,%
multicol,%
paralist,%
pgf,%
pgfplots,%
pgfplotstable,%
pgfpages,%
proof,%
psfrag,%
%subfigure,%
tcolorbox,%
tfrupee,%
tikz,%
times,%
%ulem,%
url,%
xcolor,%
mathpazo,%
}

\definecolor{shadecolor}{gray}{.95}%{rgb}{1,0,0}
\usepackage[margin=1in,top=0.75in]{geometry}
\usepackage[mathscr]{eucal}
\usepgflibrary{shapes}
\usepgfplotslibrary{fillbetween}
\usetikzlibrary{%
arrows,%
backgrounds,%
chains,%
decorations.pathmorphing,% /pgf/decoration/random steps | erste Graphik
decorations.text,% 
matrix,%
positioning,% wg. " of "
fit,%
patterns,%
petri,%
plotmarks,%
scopes,%
shadows,%
shapes.misc,% wg. rounded rectangle
shapes.arrows,%
shapes.callouts,%
shapes%
}

%\pgfplotsset{compat=newest} %<------ Here
\pgfplotsset{compat=1.11} %<------ Or use this one

\theoremstyle{plain}
\newtheorem{thm}{Theorem}[section]
\newtheorem{lem}[thm]{Lemma}
\newtheorem{prop}[thm]{Proposition}
\newtheorem{cor}[thm]{Corollary}
\newtheorem{clm}[thm]{Claim}

\theoremstyle{definition}
\newtheorem{defn}[thm]{Definition}
\newtheorem{conj}[thm]{Conjecture}
\newtheorem{exmp}[thm]{Example}
\newtheorem{axiom}[thm]{Axiom}
\newtheorem{assum}[thm]{Assumption}
\newtheorem{exerc}[thm]{Exercise}
\newtheorem*{defin}{Definition}

\theoremstyle{remark}
\newtheorem{rem}{Remark}
\newtheorem{note}{Note}

\newcommand{\iid}{\emph{i.i.d.}\ }
\newcommand{\Cov}{\operatorname{Cov}}
%\newcommand{\det}{\operatorname{det}}
%\newcommand{\Real}{\mathbb{R}}
\newcommand{\tr}{\operatorname{tr}}
\newcommand{\Var}{\operatorname{Var}}
\newcommand{\cov}{\operatorname{cov}}

%\renewcommand{\proof}[1]{\begin{proof}#1\end{proof}}
\newcommand{\EQ}[1]{\begin{equation*}#1\end{equation*}}
\newcommand{\EQN}[1]{\begin{equation}#1\end{equation}}
\newcommand{\eq}[1]{\begin{align*}#1\end{align*}}
\newcommand{\meq}[2]{\begin{xalignat*}{#1}#2\end{xalignat*}}
\newcommand{\norm}[1]{\left\lVert#1\right\rVert}
\newcommand{\inner}[1]{\left\langle #1\right\rangle}
\newcommand{\abs}[1]{\left\lvert#1\right\rvert}
\newcommand{\expect}[1]{\mathbb{E}\left[{#1}\right]}
\newcommand{\prob}[1]{\mathbb{P}\left[{#1}\right]}
\newcommand{\given}{\; \big\vert \;} 
\newcommand{\set}[1]{\left\{#1\right\}} 
\newcommand{\indicator}[1]{1_{\set{#1}}} 
\newcommand{\Ind}[1]{\mathbbm{1}_{#1}} 
\newcommand{\SetIn}[1]{\mathbbm{1}_{\set{#1}}} 
\newcommand{\red}[1]{\textcolor{red}{#1}} 
\newcommand{\floor}[1]{\left\lfloor#1\right\rfloor}
\newcommand{\ceil}[1]{\left\lceil#1\right\rceil}


\newcommand{\C}{\mathbb{C}}
\newcommand{\D}{\mathbb{D}}
\newcommand{\E}{\mathbb{E}}
\newcommand{\F}{\mathbb{F}}
\newcommand{\N}{\mathbb{N}}
\renewcommand{\P}{\mathbb{P}}
\newcommand{\Q}{\mathbb{Q}}
\newcommand{\R}{\mathbb{R}}
\newcommand{\Z}{\mathbb{Z}}

\newcommand{\bA}{\mathbf{A}}
\newcommand{\bI}{\mathbf{I}}
\newcommand{\bU}{\mathbf{1}}
\newcommand{\bx}{\mathbf{x}}
\newcommand{\bX}{\mathbf{X}}
\newcommand{\bZ}{\mathbf{Z}}


\newcommand{\cA}{\mathcal{A}}
\newcommand{\cB}{\mathcal{B}}
\newcommand{\cC}{\mathcal{C}}
\newcommand{\cF}{\mathcal{F}}
\newcommand{\cG}{\mathcal{G}}
\newcommand{\cM}{\mathcal{M}}
\newcommand{\cN}{\mathcal{N}}
\newcommand{\cP}{\mathcal{P}}
\newcommand{\cT}{\mathcal{T}}
\newcommand{\cX}{\mathcal{X}}
\newcommand{\cY}{\mathcal{Y}}

\newcommand{\sA}{\mathscr{A}}
\newcommand{\sB}{\mathscr{B}}
\newcommand{\sC}{\mathscr{C}}
\newcommand{\sE}{\mathscr{E}}
\newcommand{\sF}{\mathscr{F}}
\newcommand{\sG}{\mathscr{G}}
\newcommand{\sH}{\mathscr{H}}
\newcommand{\sL}{\mathscr{L}}
\newcommand{\sS}{\mathscr{S}}
\newcommand{\sT}{\mathscr{T}}
\newcommand{\sX}{\mathscr{X}}
\newcommand{\sY}{\mathscr{Y}}
\newcommand{\sZ}{\mathscr{Z}}

\newcommand{\tS}{\tilde{S}}
% Debug
\newcommand{\todo}[1]{\begin{color}{blue}{{\bf~[TODO:~#1]}}\end{color}}

% a few handy macros

\renewcommand{\le}{\leqslant}
\renewcommand{\ge}{\geqslant}
\newcommand\matlab{{\sc matlab}}
\newcommand{\goto}{\rightarrow}
\newcommand{\bigo}{{\mathcal O}}
%\newcommand{\half}{\frac{1}{2}}
%\newcommand\implies{\quad\Longrightarrow\quad}
\newcommand\reals{{{\rm l} \kern -.15em {\rm R} }}
\newcommand\complex{{\raisebox{.043ex}{\rule{0.07em}{1.56ex}} \hskip -.35em {\rm C}}}


% macros for matrices/vectors:

% matrix environment for vectors or matrices where elements are centered
\newenvironment{mat}{\left[\begin{array}{ccccccccccccccc}}{\end{array}\right]}
\newcommand\bcm{\begin{mat}}
\newcommand\ecm{\end{mat}}

% matrix environment for vectors or matrices where elements are right justifvied
\newenvironment{rmat}{\left[\begin{array}{rrrrrrrrrrrrr}}{\end{array}\right]}
\newcommand\brm{\begin{rmat}}
\newcommand\erm{\end{rmat}}

% for left brace and a set of choices
%\newenvironment{choices}{\left\{ \begin{array}{ll}}{\end{array}\right.}
\newcommand\when{&\text{if~}}
\newcommand\otherwise{&\text{otherwise}}
% sample usage:
%\delta_{ij} = \begin{choices} 1 \when i=j, \\ 0 \otherwise \end{choices}


% for labeling and referencing equations:
\newcommand{\eql}{\begin{equation}\label}
\newcommand{\eqn}[1]{(\ref{#1})}
% can then do
%\eql{eqnlabel}
%...
%\end{equation}
% and refer to it as equation \eqn{eqnlabel}.
% some useful macros for finite difference methods:
\newcommand\unp{U^{n+1}}
\newcommand\unm{U^{n-1}}
% for chemical reactions:
\newcommand{\react}[1]{\stackrel{K_{#1}}{\rightarrow}}
\newcommand{\reactb}[2]{\stackrel{K_{#1}}{~\stackrel{\rightleftharpoons}
 {\scriptstyle K_{#2}}}~}
\makeatletter
\def\th@plain{%
\thm@notefont{}% same as heading font
\itshape % body font
}
\def\th@definition{%
\thm@notefont{}% same as heading font
\normalfont % body font
}
\makeatother
\date{}



\usepackage{%
	amsfonts,%
	amsmath,%	
	amssymb,%
	amsthm,%
	babel,%
	bbm,%
	%biblatex,%
	caption,%
	centernot,%
	color,%
	enumerate,%
	%enumitem,%
	epsfig,%
	epstopdf,%
	etex,%
	fancybox,%
	framed,%
	fullpage,%
	%geometry,%
	graphicx,%
	hyperref,%
	latexsym,%
	mathptmx,%
	mathtools,%
	multicol,%
	paralist,%
	pgf,%
	pgfplots,%
	pgfplotstable,%
	pgfpages,%
	proof,%
	psfrag,%
	%subfigure,%	
	tcolorbox,%
	tfrupee,%
	tikz,%
	times,%
	%ulem,%
	url,%
	xcolor,%
	mathpazo,%
}

\definecolor{shadecolor}{gray}{.95}%{rgb}{1,0,0}
\usepackage[margin=1in,top=0.75in]{geometry}
\usepackage[mathscr]{eucal}
\usepgflibrary{shapes}
\usepgfplotslibrary{fillbetween}
\usetikzlibrary{%
  arrows,%
  backgrounds,%
  chains,%
  decorations.pathmorphing,% /pgf/decoration/random steps | erste Graphik
  decorations.text,% 
  matrix,%
  positioning,% wg. " of "
  fit,%
  patterns,%
  petri,%
  plotmarks,%
  scopes,%
  shadows,%
  shapes.misc,% wg. rounded rectangle
  shapes.arrows,%
  shapes.callouts,%
  shapes%
}

%\pgfplotsset{compat=newest} %<------ Here
\pgfplotsset{compat=1.11} %<------ Or use this one

\theoremstyle{plain}
\newtheorem{thm}{Theorem}[section]
\newtheorem{lem}[thm]{Lemma}
\newtheorem{prop}[thm]{Proposition}
\newtheorem{cor}[thm]{Corollary}
\newtheorem{clm}[thm]{Claim}

\theoremstyle{definition}
\newtheorem{defn}[thm]{Definition}
\newtheorem{conj}[thm]{Conjecture}
\newtheorem{exmp}[thm]{Example}
\newtheorem{axiom}[thm]{Axiom}
\newtheorem{assum}[thm]{Assumption}
\newtheorem{exerc}[thm]{Exercise}
\newtheorem*{defin}{Definition}

\theoremstyle{remark}
\newtheorem{rem}{Remark}
\newtheorem{note}{Note}

\newcommand{\iid}{\emph{i.i.d.}\ }
\newcommand{\Cov}{\operatorname{Cov}}
%\newcommand{\det}{\operatorname{det}}
%\newcommand{\Real}{\mathbb{R}}
\newcommand{\tr}{\operatorname{tr}}
\newcommand{\Var}{\operatorname{Var}}
\newcommand{\cov}{\operatorname{cov}}

%\renewcommand{\proof}[1]{\begin{proof}#1\end{proof}}
\newcommand{\EQ}[1]{\begin{equation*}#1\end{equation*}}
\newcommand{\EQN}[1]{\begin{equation}#1\end{equation}}
\newcommand{\eq}[1]{\begin{align*}#1\end{align*}}
\newcommand{\meq}[2]{\begin{xalignat*}{#1}#2\end{xalignat*}}
\newcommand{\norm}[1]{\left\lVert#1\right\rVert}
\newcommand{\inner}[1]{\left\langle #1\right\rangle}
\newcommand{\abs}[1]{\left\lvert#1\right\rvert}
\newcommand{\expect}[1]{\mathbb{E}\left[{#1}\right]}
\newcommand{\prob}[1]{\mathbb{P}\left[{#1}\right]}
\newcommand{\given}{\; \big\vert \;} 
\newcommand{\set}[1]{\left\{#1\right\}} 
\newcommand{\indicator}[1]{1_{\set{#1}}} 
\newcommand{\Ind}[1]{\mathbbm{1}_{#1}} 
\newcommand{\SetIn}[1]{\mathbbm{1}_{\set{#1}}} 
\newcommand{\red}[1]{\textcolor{red}{#1}} 
\newcommand{\floor}[1]{\left\lfloor#1\right\rfloor}
\newcommand{\ceil}[1]{\left\lceil#1\right\rceil}


\newcommand{\C}{\mathbb{C}}
\newcommand{\D}{\mathbb{D}}
\newcommand{\E}{\mathbb{E}}
\newcommand{\F}{\mathbb{F}}
\newcommand{\N}{\mathbb{N}}
\renewcommand{\P}{\mathbb{P}}
\newcommand{\Q}{\mathbb{Q}}
\newcommand{\R}{\mathbb{R}}
\newcommand{\Z}{\mathbb{Z}}

\newcommand{\bA}{\mathbf{A}}
\newcommand{\bI}{\mathbf{I}}
\newcommand{\bU}{\mathbf{1}}
\newcommand{\bx}{\mathbf{x}}
\newcommand{\bX}{\mathbf{X}}
\newcommand{\bZ}{\mathbf{Z}}


\newcommand{\cA}{\mathcal{A}}
\newcommand{\cB}{\mathcal{B}}
\newcommand{\cC}{\mathcal{C}}
\newcommand{\cF}{\mathcal{F}}
\newcommand{\cG}{\mathcal{G}}
\newcommand{\cM}{\mathcal{M}}
\newcommand{\cN}{\mathcal{N}}
\newcommand{\cP}{\mathcal{P}}
\newcommand{\cT}{\mathcal{T}}
\newcommand{\cX}{\mathcal{X}}
\newcommand{\cY}{\mathcal{Y}}

\newcommand{\sA}{\mathscr{A}}
\newcommand{\sB}{\mathscr{B}}
\newcommand{\sC}{\mathscr{C}}
\newcommand{\sE}{\mathscr{E}}
\newcommand{\sF}{\mathscr{F}}
\newcommand{\sG}{\mathscr{G}}
\newcommand{\sH}{\mathscr{H}}
\newcommand{\sL}{\mathscr{L}}
\newcommand{\sS}{\mathscr{S}}
\newcommand{\sT}{\mathscr{T}}
\newcommand{\sX}{\mathscr{X}}
\newcommand{\sY}{\mathscr{Y}}
\newcommand{\sZ}{\mathscr{Z}}

\newcommand{\tS}{\tilde{S}}
% Debug
\newcommand{\todo}[1]{\begin{color}{blue}{{\bf~[TODO:~#1]}}\end{color}}

% a few handy macros

\renewcommand{\le}{\leqslant}
\renewcommand{\ge}{\geqslant}
\newcommand\matlab{{\sc matlab}}
\newcommand{\goto}{\rightarrow}
\newcommand{\bigo}{{\mathcal O}}
%\newcommand{\half}{\frac{1}{2}}
%\newcommand\implies{\quad\Longrightarrow\quad}
\newcommand\reals{{{\rm l} \kern -.15em {\rm R} }}
\newcommand\complex{{\raisebox{.043ex}{\rule{0.07em}{1.56ex}} \hskip -.35em {\rm C}}}


% macros for matrices/vectors:

% matrix environment for vectors or matrices where elements are centered
\newenvironment{mat}{\left[\begin{array}{ccccccccccccccc}}{\end{array}\right]}
\newcommand\bcm{\begin{mat}}
\newcommand\ecm{\end{mat}}

% matrix environment for vectors or matrices where elements are right justifvied
\newenvironment{rmat}{\left[\begin{array}{rrrrrrrrrrrrr}}{\end{array}\right]}
\newcommand\brm{\begin{rmat}}
\newcommand\erm{\end{rmat}}

% for left brace and a set of choices
%\newenvironment{choices}{\left\{ \begin{array}{ll}}{\end{array}\right.}
\newcommand\when{&\text{if~}}
\newcommand\otherwise{&\text{otherwise}}
% sample usage:
%  \delta_{ij} = \begin{choices} 1 \when i=j, \\ 0 \otherwise \end{choices}


% for labeling and referencing equations:
\newcommand{\eql}{\begin{equation}\label}
\newcommand{\eqn}[1]{(\ref{#1})}
% can then do
%  \eql{eqnlabel}
%  ...
%  \end{equation}
% and refer to it as equation \eqn{eqnlabel}.  
% some useful macros for finite difference methods:
\newcommand\unp{U^{n+1}}
\newcommand\unm{U^{n-1}}
% for chemical reactions:
\newcommand{\react}[1]{\stackrel{K_{#1}}{\rightarrow}}
\newcommand{\reactb}[2]{\stackrel{K_{#1}}{~\stackrel{\rightleftharpoons}
   {\scriptstyle K_{#2}}}~}
\makeatletter
\def\th@plain{%
  \thm@notefont{}% same as heading font
  \itshape % body font
}
\def\th@definition{%
  \thm@notefont{}% same as heading font
  \normalfont % body font
}
\makeatother
\date{}



\title{Lecture-02: Probability Function}
\author{}

\begin{document}
\maketitle

\section{Indicator Functions}
\begin{defn}[Indicator functions]
Consider a random experiment over an outcome space $\Omega$ and an event space $\sF$. 
The indicator of an event $B \in \sF$ is denoted by 
\EQ{
\Ind{B}(\omega) = 
\begin{cases}
1, & \omega \in B,\\
0, & \omega \notin B.
\end{cases}
}
\end{defn}
\begin{shaded}
\begin{exmp} 
Consider a roll of dice that has an outcome space $\Omega = [6]$ and event space $\sF = \cP(\Omega)$. 
For an event $O \triangleq \set{1,3,5}$ that represents odd outcomes for dice roll, we observe that $\Ind{O}(1) = 1$ and $\Ind{O}(2) = 0$. 
\end{exmp}
\begin{exmp}
Consider $N$ trials of a random experiment over outcome space $\Omega$ and the event space $\sF$.  
Let $\omega_n \in \Omega$ denote the outcome of the experiment of the $n$th trial. 
For each event $B \in \sF$, we define an indicator function  
\EQ{
\Ind{B}(\omega_n) = 
\begin{cases}
1, & \omega_n \in B,\\
0, & \omega_n \notin B.
\end{cases}
}
For any event $B \in \sF$, 
the number of times the event $B$ occurs in $N$ trials is denoted by $N(B) = \sum_{n=1}^N\Ind{B}(\omega_n)$. 
We denote the relative frequency of an event $A$ in $N$ trials by $\frac{N(B)}{N}$. 
\end{exmp}
\end{shaded}
\begin{defn}[Disjoint events]
Let $(\Omega,\sF)$ be a pair of sample and event space, 
and a sequence of events $A \in \sF^\N$. 
The sequence is \textbf{mutually disjoint} if $A_n\cap A_m = \emptyset$ for any $m\neq n \in \N$. 
If $\cup_{n \in \N}A_n = \Omega$, then $A$ is a \textbf{partition} of sample space $\Omega$.   
\end{defn}
\begin{tcolorbox}
\begin{exerc}
Consider the sample space $\Omega = \set{H,T}^\N$ and event space $\sF = \sigma(A_n: n \in \N)$ where $A_n \triangleq \set{\omega \in \Omega: \omega_i = H \text{ for some } i \in [n]}$. 
Construct a sequence of non-trivial disjoint events. 
\end{exerc}
\begin{exerc}
Let $(\Omega,\sF)$ be a pair of sample and event space. 
For any sequence of events $A \in \sF^\N$, show that
\begin{enumerate}
\item $\Ind{\cap_{n\in\N}A_n}(\omega) = \prod_{n \in \N}\Ind{A_n}(\omega)$, 
\item $\Ind{\cup_{n \in \N}A_n}(\omega) = \sum_{n \in \N}\Ind{A_n}(\omega)$ if the sequence $A$ is mutually disjoint.   
\end{enumerate} 
\end{exerc}
\end{tcolorbox}
\begin{shaded}
\begin{exmp}
We observe the following properties of the relative frequency.
\begin{enumerate}
\item For all events $B \in \sF$, we have $0 \le \frac{N(B)}{N} \le 1$. 
This follows from the fact that $0 \le N(B) \le N$ for any event $B \in \sF$. 
\item Let $A\in\sF^\N$ be a sequence of mutually disjoint events,   
%Suppose $A_i \in \sF$ for all $i \in \N$ such that $A_i \cap A_j = \emptyset$ for all $i \neq j$, 
then $\frac{N(\cup_{i \in \N}A_i)}{N} =  \sum_{i\in \N}\frac{N(A_i)}{N}$. 
This follows from the fact that for mutually disjoint event sequence $A$, we have 
\EQ{
\Ind{\cup_{i \in \N}A_i}(\omega_n) = \sum_{i \in \N}\Ind{A_i}(\omega_n).
}
\item For the certain event $\Omega$, we have $\frac{N(\Omega)}{N} = 1$. 
This follows from the fact that $N(\Omega) = N$. 
\end{enumerate}
Since the relative frequency is positive and bounded, it may converge to a real number as $N$ grows very large, 
and the limit $\lim_{N \to \infty}\frac{N(B)}{N}$ may exist. 
\end{exmp}
\end{shaded} 
\section{Probability axioms}
Inspired by the relative frequency, we list the following axioms for a probability function $P: \sF \to [0,1]$. 
\begin{axiom}[Axioms of probability]
We define a probability measure on sample space $\Omega$ and event space $\sF$ by a function $P: \sF \to [0,1]$ which satisfies the following axioms. 
\begin{enumerate}
\item  [\textbf{Non-negativity:}] For all events $B \in \sF$, we have $P(B) \ge 0$. 
%\item  [\textbf{Finite additivity}] For any two sets of mutually disjoint events $A, B \in \sF$ such that $A \cap B = \emptyset$, we have $P(A \cup B) = P(A) + P(B)$. 
\item  [\textbf{$\sigma$-additivity:}] For an infinite sequence $A \in \sF^\N$ of mutually disjoint events,  
%$A_i \in \sF$ for all $i \in \N$ such that $A_i \cap A_j = \emptyset$ for all $i \neq j$, 
we have $P(\cup_{i \in \N}A_i) = \sum_{i \in \N}P(A_i)$. 
\item [\textbf{Certainty:}] $P(\Omega) = 1$. 
\end{enumerate} 
\end{axiom}
\begin{defn}[Probability space] 
A sample space $\Omega$, an event space $\sF\subseteq \cP(\Omega)$, 
and a probability measure $P: \sF \to [0,1]$, 
together define a probability space $(\Omega,\sF,P)$.
\end{defn}
\section{Properties of Probability}
\begin{thm}
For any probability space $(\Omega, \sF, P)$, 
we have the following properties of probability measure. 
\begin{enumerate}
\item [\textbf{impossibility:}] $P(\emptyset) = 0$.
\item [\textbf{finite additivity:}] For mutually disjoint events $A\in \sF^n$, 
%such that $A_i \cap A_j =\emptyset$ for $i \neq j$, 
we have $P(\cup_{i=1}^nA_i) = \sum_{i=1}^nP(A_i)$. 
\item [\textbf{monotonicity:}] If events $A, B \in \sF$ such that $A \subseteq B$, then $P(A) \le P(B)$. 
\item [\textbf{inclusion-exclusion:}] For any events $A, B \in \sF$, we have $P(A \cup B) = P(A) + P(B) - P(A \cap B)$.
\item [\textbf{continuity:}] For a sequence of events $A \in \sF^\N)$ such that $\lim_nA_n$ exists, we have $P(\lim_{n \to \infty}A_n) = \lim_{n \to \infty}P(A_n)$. 
\end{enumerate}
\end{thm}
\begin{proof}
We consider the probability space $(\Omega, \sF, P)$.
\begin{enumerate}
\item We take disjoint events $E\in\sF^\N$ where $E_1 = \Omega$ and $E_i = \emptyset$ for $i \ge 2$. 
It follows that $\cup_{i \in \N}E_i = \Omega$ and $E$ is a collection of mutually disjoint events.  
From the countable additivity axiom of probability, it follows that
\EQ{
P(\Omega) = P(\Omega) + \sum_{i \ge 2}P(E_i). 
}
Since $P(E_i) \ge 0$, it implies that $P(\emptyset) = 0$. 
\item 
We see that finite additivity follows from the countable additivity. 
We consider disjoint events $A_1, \dots, A_n$, and take $A_{i} = \emptyset$ for all $i > n$. 
It follows that the sequence of sets $A\in\sF^\N$ is mutually disjoint, and since $P(\emptyset) = 0$, it follows that  
\EQ{
P(\cup_{i=1}^nA_i) = P(\cup_{i \in \N}A_i) = \sum_{i=1}^nP(A_i) + \sum_{i > n}P(\emptyset) =  \sum_{i=1}^nP(A_i).
}
\item For events $A, B \in \sF$ such that $A \subseteq B$, we can take disjoint events $E_1 = A$ and $E_2 = B\setminus A$. 
From closure under complements and intersection, it follows that $E_2 \in \sF$. 
From non-negativity of probability, we have $P(E_2) \ge 0$. 
Finally, the result follows from finite additivity of disjoint events 
\EQ{
P(B) = P(E_1 \cup E_2)= P(E_1) + P(E_2) \ge P(A).
}
\item For any two events $A, B \in \sF$, we can write the following events as disjoint unions  
\meq{3}{
&A = (A\setminus B) \cup (A\cap B), &&B = (B\setminus A) \cup (A\cap B), && A\cup B = (A\setminus B)\cup(A \cap B)\cup(B \setminus A).
}
The result follows from the finite additivity of probability of disjoint events. 

\item To show the continuity of probability in events, we first need to understand the limits of events. 
We show the continuity of probability in the next section. 
%We show the continuity for non-decreasing and non-increasing sequence of sets.\\
\end{enumerate}
\end{proof}
\begin{shaded}
\begin{exmp}
Consider a single coin tosse with the sample space $\Omega = \set{H,T}$, an event space $\sF = \cP(\Omega)$. 
The probability measure $P:\sF\to[0,1]$ is defined as 
\begin{xalignat*}{4}
&P(\emptyset) = 0,&
&P(\set{H}) = p&
&P(\set{T}) = 1-p&
&P(\set{H,T}) = 1.
\end{xalignat*}
Can you verify that $P$ is a probability function?
\end{exmp}
\begin{exmp}
Consider $N$ coin tosses with the sample space $\Omega = \set{H,T}^{[N]}$, an event space $\sF = \cP(\Omega)$. 
For each outcome $\omega \in \Omega$, we define the number of heads as $N_H(\omega) \triangleq \sum_{n\in[N]}\Ind{\set{H}}(\omega_n)$ and the number of tails as $N_T(\omega) \triangleq N - N_H(\omega)$. 
The probability measure $P:\sF\to[0,1]$ is defined for each event $A \in \sF$ as 
\begin{equation*}
P(A) \triangleq \sum_{\omega \in A}p^{N_H(\omega)}(1-p)^{N_T(\omega)}.
\end{equation*}
Can you verify that $P$ is a probability function?
\end{exmp}
\end{shaded}
\section{Limits of Sets}
\begin{defn}[Limits of monotonic sets]
For a sequence of non-decreasing sets $(A_n: n \in \N)$, we can define the limit as 
\EQ{
\lim_{n \to \infty}A_n \triangleq \cup_{n \in \N}A_n.
}
Similarly, for a sequence of non-increasing sets $(A_n: n \in \N)$, we can define the limit as 
\EQ{
\lim_{n \to \infty}A_n \triangleq \cap_{n \in \N}A_n.
}
\end{defn}
\begin{shaded}
\begin{exmp}[Monotone sets]
Consider a monotonically increasing sequence $a \in \R^\N$ defined as $a_n \triangleq -\frac{1}{n}$ for all  $n \in \N$,  
which converges to the limit $0$.  
We consider monotone sequence of sets $A, B \in \sB(\R)^\N$ define as $A_n \triangleq [-2, -\frac{1}{n}]$ and $B_n \triangleq [-2, \frac{1}{n}]$  for all $n \in \N$, which are monotonically increasing and decreasing respectively. 
We can verify the following limits 
\meq{2}{
&\lim_n A_n = \cup_{n \in \N}A_n = [-2, 0), &&\lim_nB_n = \cap_{n \in \N}B_n = [-2, 0].
}
\end{exmp}
\end{shaded}
%We now generalize this definition to general sequence of sets.
\begin{defn}[Limits of sets] 
For a sequence of sets $(A_n: n \in \N)$, we can define the limit superior and limit inferior of this sequence of sets as 
\meq{2}{
&\limsup_{n \to \infty}A_n \triangleq \cap_{n \in \N}\cup_{m \ge n}A_m = \lim_{n \to \infty}\cup_{m \ge n}A_m,
&&\liminf_{n\to\infty}A_n\triangleq \cup_{n \in \N}\cap_{m \ge n}A_m =  \lim_{n \to \infty}\cap_{m \ge n}A_m.
}
%In general, $\liminf_{n \to \infty}A_n \subseteq \limsup_{n \to \infty}A_n$. 
%If the limit superior and limit inferior of any sequence of sets are equal, then the limit is defined as 
%\EQ{
%\lim_{n \to \infty}A_n = \limsup_{n \to \infty}A_n = \liminf_{n \to \infty}A_n. 
%}
\end{defn}
\begin{lem} 
For a sequence of sets $(A_n: n \in \N)$, we have $\liminf_{n \to \infty}A_n \subseteq \limsup_{n \to \infty}A_n$.  
\end{lem}
\begin{proof}
For each $n \in \N$, we define $E_n \triangleq \cup_{m \ge n}A_m$ and $F_n \triangleq \cap_{m \ge n}A_m$. 
%We see that $F_n \subseteq A_m$ for all $m \ge n$.  
Consider a fixed $n$, then $F_1, F_2, \dots, F_n \subseteq A_n$, and $F_m \subseteq A_m$ for all $m \ge n$. 
Therefore, we can write $\cup_{n \in \N}F_n \subseteq \cup_{m \ge n}A_m$ for each $n \in \N$, 
and hence the result follows. 
%Let $x \in \liminf_n A_n$, then $x \in F_{n_0}$ for some $n_0 \in \N$.
%That is, there exists an $n_0 \in \N$, such that $x \in A_m$ for all $m \ge n_0$. 
%In particular, it means $x \in \cup_{m \ge n}A_m$ for each $n \in \N$, and hence $x \in \limsup_n A_n$. 
\end{proof}
\begin{defn}
If the limit superior and limit inferior of any sequence of sets $(A_n: n \in \N)$ are equal, then the sequence of sets has a limit $A_\infty$, which is defined as 
\EQ{
A_\infty \triangleq \lim_{n \to \infty}A_n = \limsup_{n \to \infty}A_n = \liminf_{n \to \infty}A_n. 
}
\end{defn}
\begin{shaded}
\begin{exmp}[Sequence of sets with different limits]
We consider sequence of sets $(A_n = [-2, (-1)^n + \frac{1}{n}]: n \in \N)$. 
It follows that  $F_n = \cap_{m \ge n}A_m = [-2, -1]$ and 
\EQ{
E_n = \cup_{m \ge n}A_m = 
\begin{cases}
[-2, 1+ \frac{1}{n+1}], & n \text{ odd},\\
[-2, 1 + \frac{1}{n}], &n \text{ even}.
\end{cases}
}
We can verify the following limits 
\meq{2}{
&\liminf_n A_n = \cup_{n \in \N}F_n = [-2, -1], &&\limsup_nA_n = \cap_{n \in \N}E_n = [-2, 1].
}
\end{exmp}
\end{shaded}
\subsection{Proof of continuity of probability}
\textbf{Continuity for increasing sets.} 
Let $(A_n \in \sF: n \in \N)$ be a non-decreasing sequence of events, 
%such that $A_n \subseteq A_{n+1}$ for all $n \in \N$. 
%Then, we have 
then $\lim_{n \to \infty}A_n = \cup_{n\in \N}A_n$. 
This implies that $(P(A_n): n \in \N)$ is a non-negative non-decreasing bounded sequence, 
and hence has a limit. 
It remains to show that $\lim_{n \to \infty}P(A_n) = P(\cup_{n \in \N}A_n)$. 
To this end, we observe that $(A_1, A_2\setminus A_1, \dots, A_n\setminus A_{n-1})$ is a partition of the event $A_n$, and $P(A_i\setminus A_{i-1}) = P(A_i) - P(A_{i-1})$ for each $i \in \N$. 
From finite additivity of $P$ for mutually disjoint events, we can write for each $n \in \N$
\EQ{
P(A_n) = P(A_1) + \sum_{i=1}^{n-1}P(A_{i+1}\setminus A_i) = P(A_1) + \sum_{i=1}^{n-1}(P(A_{i+1}) -P(A_i)).
}
From $\sigma$-additivity of $P$ for sequence of mutually disjoint events, 
we can write for $\lim_{n \to \infty}A_n = \cup_{n \in \N}A_n$, 
\EQ{
P(\cup_{n \in \N}A_n) = P(A_1) + \sum_{i \in \N}(P(A_{i+1}) -P(A_i)) = P(A_1) + \lim_{n \to \infty}\sum_{i=1}^{n-1}(P(A_{i+1}) -P(A_i)) = \lim_{n \to \infty}P(A_n).
}

\textbf{Continuity for decreasing sets.} Similarly, for a non-increasing sequence of sets $(B_n \in \sF: n \in \N)$, 
we can find the non-decreasing sequence of sets $(B^c_n \in \sF: n \in \N)$. 
By the first part, we have 
\EQ{
P(\lim_{n \to \infty}B_n) = P(\cap_{n \in \N}B_n) = 1 - P(\cup_{n \in \N}B^c_n)  = 1 - P(\lim_{n \to \infty}B_n^c) = 1 - \lim_{n\to \infty}P(B_n^c) = \lim_{n \to \infty}P(B_n).
}

\textbf{Continuity for general sequence of sets.} 
We can similarly prove the general result for a sequence of sets $(A_n \in \sF: n \in \N)$ such that the limits $\lim_{n}A_n$ 
%and $\lim_nP(A_n)$ exist. 
exists. 
We can define non-increasing sequences of sets $(E_n = \cup_{m \ge n}A_m \in \sF: n \in \N)$ and non-decreasing sequences of sets $(F_n = \cap_{m \ge n}A_m \in \sF: n \in \N)$. 
From the continuity of probability for the monotonic sets, we have 
\meq{2}{
&P(\limsup_nA_n) = P(\cap_{n\in\N}E_n) = \lim_{n\to \infty}P(E_n),&&P(\liminf_nA_n) = P(\cup_{n\in\N}F_n) = \lim_{n\to \infty}P(F_n).
}
From the definition of two sequences of sets, we obtain 
\meq{2}{
&P(E_n) \ge \sup_{m \ge n}P(A_m), && P(F_n) \le \inf_{m \ge n}P(A_m).
} 
Therefore taking limsup and liminf, we obtain 
\EQ{
\lim_{n\to\infty}P(E_n) = \limsup_{n \in \N} P(A_n) \ge \inf_{n \in \N}\sup_{m \ge n}P(A_m) \ge  \sup_{n \in \N}\inf_{m \ge n}P(A_m) \ge \liminf_{n \in \N}P(A_n)
= \lim_{n\to\infty}P(F_n).
} 
Since $P(\lim_n A_n) = \lim_nP(E_n) = \lim_{n}P(F_n)$ exists, the result follows.

\end{document}
